% ============================================================
% sections/intro.tex  はじめに（初稿 2026-10-26）
% 数字は入れない（結果の予告は向きだけ）。引用キーは references.bib と対応。
% 文献の確かめの状況は執筆メモ §11 を参照。
% ============================================================
\section{はじめに}\label{sec:intro}

他者の振る舞いを予想し，その理由を理解する能力は，社会的な生活の基盤である．
この能力の説明として，二つの立場が長く論じられてきた．
理論説（theory theory）は，他者の心的状態を，信念や欲求についての素朴な理論に基づく推論として説明する\cite{gopnik1994theory}．
シミュレーション説（simulation theory）は，他者の立場に自分を置き，自分の心の働きを用いて他者の心的状態を推し量ると考える\cite{gordon1986folk,goldman2006simulating}．
シミュレーション説の神経科学的な支持として，自分が行為するときと他者の行為を観察するときに共通して活動する神経活動が挙げられてきた\cite{gallese1998mirror,rizzolatti2004mirror}．

しかし，どちらの立場も，「自己」と「他者」という枠組みが最初から与えられていることを前提にしている．
Noguchi ら\cite{noguchi2022superposition}はこの前提を問い，自他の区別を与えずに，予測学習だけから自他の表現が生まれる仕組みとして重ね合わせ機構（superposition mechanism）を提案した．
この機構では，一つの感覚入力を二つの経路で符号化し，それらを重みを共有した一つのモジュールで処理する．
文献\cite{noguchi2022superposition}の新しさは，自分の経験が他者の予測に役立つこと自体にあるのではない．
他者の動きを自分のどの経験に当てはめればよいかという自他の対応を一切与えずに，自分の次の視界を予測する学習だけから，自他の区別，他者の位置，他者の視点，他者の運動が同時に生まれた点にある．
すなわち，共有モジュールの中に自分と他者の位置が同じ形で表され（社会的場所細胞\cite{danjo2018spatial,omer2018social}に似た表現），他者の視点から見た視界を取り出すことができ，他者の運動を自分の運動と同じ形で生成するモジュールが学習された．
これは，一つのモジュールを「誰にでも」使うことで，自分を理解するのと同じ仕方で他者を理解しうることを示している．

一方，文献\cite{noguchi2022superposition}自身が述べているように，この機構が扱ったのは他者の位置・視点・運動までであり，他者の考えや欲求のような内面には届いていない．
また文献\cite{noguchi2022superposition}は，現実の発達では自分についての経験と他者についての経験の量に非対称があり，それが自他の理解の形成に影響しうることを今後の課題に挙げている．
本研究はこの二つの課題を結びつける．
すなわち，自分の価値に従って動いた経験の量を操作し，その経験が重ね合わせ機構を通して他者の理解にどう及ぶかを調べる．
本研究は，文献\cite{noguchi2022superposition}が示した自他の対応を土台として引き継ぐ．
この土台は本研究の新しさではない．本研究では，自他の対応を担う視覚エンコーダを，文献\cite{noguchi2022superposition}の設定で学習した重みに固定して与えている（\ref{sec:network}節）．
本研究が問うのは，その対応に乗って自分の価値に従う見込みが他者に運ばれるとき，何が得られ，何を代償とし，どこに限界があるかである．

シミュレーションによる他者理解には，できることとできないことがあるはずである．
自分の心の働きを他者に当てはめれば，他者を自分と同じく価値に従って動く者として予測できる．
同時に，自分と異なる他者に対しては，自分の性質を他者に当てはめる誤り（投影）を生みうる．
人間でも，自分の知識や好み，その時の身体の状態を他者に当てはめてしまう傾向が繰り返し報告されている\cite{ross1977false,nickerson1999how,vanboven2003social,epley2004perspective}．
人の脳では，他者の価値に基づく選択を予測するときに自分の価値づけの仕組みが流用され，他者の価値が自分と異なるときには，他者の行動の予測誤差による補正が重要になることが示唆されている\cite{suzuki2012learning}．
さらに，他者の動きを見てその好みを読み取ることは，自分の心の働きを当てはめることとは別の働きであり，シミュレーションだけで生まれるかは自明でない．
発達研究では，自分で物をつかむ経験が乳児の他者の行為の目標の理解を変えること\cite{sommerville2005action}や，14か月児が自分と違う好みを示した他者にも自分の好む物を差し出すこと\cite{repacholi1997early}が報告されている．
本研究の問いはこれらと類推的に対応する（\ref{sec:related-dev}節）．

そこで本研究は，次の主の問いを立てる．
\begin{quote}
  自他で共有された予測の仕組みは，自分の価値の経験から，他者について何を生み，何を生まないか．
\end{quote}
ここで価値は，方向ごとの快・不快として人が与える信号である（\ref{sec:env}節）．
本研究が問うのは価値そのものが生まれるかではなく，与えた価値の信号と，それに従う行動の規則性を，自他で共有された予測学習がどう利用するかである．
この問いを四つに分ける．
(RQ1) 自分の価値の経験は，予測学習だけで価値に従う行動の見込みを生むか．その強さは経験の量で変わるか．
(RQ2) その見込みは共有モジュールを通って他者側に及ぶか．
(RQ3) 他者用のモジュールを学習させると，他者の好みを自ら書き出すようになるか．
(RQ4) 見込みを用いて他者の好みの候補を入れ比べると，他者の好みを見分けられるか．
それぞれに対して，次の予測を検証した．
(H1) 価値に従う行動の見込みは，価値に従って動いた経験があるときに経験なしよりはっきり強く生まれ，経験が多いほど強い．
(H2) その見込みは共有モジュールを通して他者側に入れた快・不快に働き，共有をなくすと他者側の働きだけが消える．
(H3) 他者用のモジュールは，他者の快・不快を表すようになる．
(H4) 経験が多いとき，候補の照合により他者の好みを偶然より高く見分けられ，その見分けは他者が直前に進んだ方角では説明できない．
H1・H2 は支持され，H3 は支持されなかった．H4 は前半が支持された．後半は，一部のモデルで方角に反して好みの側を選ぶ傾向が見られたが，好みに固有の情報による判別は確証されなかった．
ただし，経験の量を変えると学習データの位置の分布も変わり，H1・H2 の経験の効果はこの変化と分けられていない（\ref{sec:limitations}節）．
加えて，予想していなかった結果として，行動の仕方が自分と違う他者に対しては，経験の最も多い条件で，本当の快・不快を与えたときの予測の誤差が大きくなり，経験の少ない条件より外れることを見いだした（経験の多い条件の4組のうち2〜3組では，何も与えない場合より悪かった）．
追加の比較では，他者が好みに従う規則を選ぶ確率を自分の学習条件にそろえると，本当の快・不快を与えたときの，何も与えない場合に対する相対的な誤差が改善した．

本研究の貢献は，既存の重ね合わせ機構を価値に関わる信号の扱いへ広げ，自分の経験から学ばれた予測の傾向が他者側に転用されること，その使い道と限界を，一つの計算モデルの上で系統的に調べた点にある．
\begin{enumerate}
  \item 見込みは他者に運ばれる．このモデルでは，自分側で学んだ重みが他者側の計算に使われることがそのために要る．自他の重みの共有だけを外すと，自分側の見込みはそのままで，他者側の働きだけが消える（ただし，この対照では他者側の快・不快を受ける重みが学習されないので，結果は構成からほぼ予想されるものである）．
  \item 自分側の見込みは，価値に従って動いた経験の量に応じて段階的に強まり，それが共有を通して他者に当てはめられる．
  \item 行動の仕方が自分と違う他者では，経験の多い条件で，本当の快・不快を与えたときの予測の誤差が経験の少ない条件より大きくなった（経験の多い条件の4組のうち2〜3組では，何も与えない場合より悪かった）．この現象は，自分の行動の傾向を他者に当てはめすぎる投影の誤りとして解釈できるが，他者の規則の確率を自分の学習条件にそろえると相対的な誤差が改善することから，学習で見ていない入力であることだけでは説明できない．ただし，学習データの位置の分布による説明は残る（\ref{sec:res-projection}節）．
  \item 運ばれた見込みは，限られた形で使える．研究者が用意した二つの好みの候補を入れ比べる照合では，経験が多いときに偶然を上回る．ただし主の基準での正解率は他者の進んだ方角だけで選ぶ基準を上回らず（副の基準でもわずかに上回るにとどまる），方角と食い違う場面で好みに従うのは一部の組に限られる．一方，他者用のモジュールが学ぶ補正は，他者の価値を直す形ではなく，他者の位置や方角という予測に役立つ手がかりを伝える形をとり，調べた4通りの構成のどれでも他者の好みを出力に書き出さない．
\end{enumerate}

以下，\ref{sec:related}章で関連研究を整理し，\ref{sec:method}章でモデルと実験方法を述べる．
\ref{sec:results}章で四つの実験の結果を示し，\ref{sec:discussion}章でシミュレーション説と発達研究との関係，および限界を論じる．
\ref{sec:conclusion}章で全体をまとめる．